\chapter{Analyse du besoin et cahier des charges}

\section{Situation observée}

La ligne peut produire 24 heures sur 24, cinq jours par semaine. Cela ne veut
pas dire que l'imprimante doit fonctionner en continu. Son utilisation dépend
des commandes et de la référence de pare-brise en cours. Pendant une période
sans demande, une imprimante arrêtée n'est pas en panne. À l'inverse, si un
verre arrive au poste et que l'imprimante est indisponible, le temps perdu doit
être enregistré.

Cette différence paraît simple, mais elle change les indicateurs. Une
disponibilité calculée sur le temps calendaire pénaliserait les périodes sans
commande. Elle produirait des alertes inutiles, puis les utilisateurs
finiraient par les ignorer.

\section{Parties prenantes}

\begin{table}[H]
\centering
\caption{Attentes des utilisateurs}
\begin{tabularx}{\textwidth}{p{0.22\textwidth}YY}
\toprule
Utilisateur & Question principale & Information attendue \\
\midrule
Opérateur &
La machine peut-elle imprimer maintenant ? &
État, niveau des fluides, défaut actif et volume récent. \\
Maintenance &
Pourquoi la machine s'arrête-t-elle et quand intervenir ? &
Tendances, heures restantes, historique, Pareto des défauts et service à prévoir. \\
Direction &
Quel est l'effet sur la production ? &
Disponibilité liée à la demande, temps d'arrêt, rendement et volume marqué. \\
\bottomrule
\end{tabularx}
\end{table}

\section{Exigences fonctionnelles}

\begin{longtable}{p{0.11\textwidth}p{0.58\textwidth}p{0.20\textwidth}}
\caption{Exigences fonctionnelles du prototype}\\
\toprule
Réf. & Exigence & Critère d'acceptation \\
\midrule
\endfirsthead
\toprule
Réf. & Exigence & Critère d'acceptation \\
\midrule
\endhead
F01 & Recevoir les mesures et événements par MQTT. &
Les trois familles de topics sont consommées. \\
F02 & Afficher l'état courant de l'imprimante. &
État mis à jour en moins d'une période de publication. \\
F03 & Suivre les niveaux d'encre et de solvant. &
Courbes, valeur courante et seuils visibles. \\
F04 & Estimer le temps avant épuisement. &
Prévision calculée à partir d'une pente récente. \\
F05 & Gérer la péremption d'un consommable. &
Date imprimée et durée après ouverture comparées. \\
F06 & Distinguer un arrêt réel d'une absence de commande. &
Le signal de demande conditionne l'indicateur. \\
F07 & Enregistrer les défauts et leur durée. &
Historique filtrable et Pareto sur 30 jours. \\
F08 & Conserver le marquage DMC et son résultat de contrôle. &
Recherche par DMC et statistiques de qualité possibles. \\
F09 & Déclencher des alertes. &
Sept règles provisionnées et un webhook testé. \\
F10 & Présenter des vues adaptées aux trois publics. &
Trois tableaux de bord distincts. \\
\bottomrule
\end{longtable}

\section{Exigences non fonctionnelles}

\begin{itemize}
  \item \textbf{Lecture seule :} aucun service du projet n'écrit vers le PLC ou
        l'imprimante.
  \item \textbf{Reproductibilité :} l'environnement complet doit démarrer avec
        Docker Compose.
  \item \textbf{Traçabilité :} chaque message possède un identifiant unique.
  \item \textbf{Résistance aux doublons :} une livraison MQTT répétée ne doit
        pas créer deux mesures ou deux événements.
  \item \textbf{Observabilité :} chaque service expose un contrôle de santé ou
        un état visible dans Docker.
  \item \textbf{Confidentialité :} les mots de passe ne figurent ni dans les
        flux Node-RED ni dans le fichier Compose.
  \item \textbf{Évolutivité raisonnable :} l'ajout d'un tag ne doit pas imposer
        une refonte du modèle de stockage.
\end{itemize}

\section{Périmètre et hypothèses}

Le périmètre s'arrête à la supervision. Les commandes de production, l'ERP, le
MES, le réglage de l'imprimante et l'écriture dans les automates sont exclus.
Le prototype utilise les hypothèses suivantes :

\begin{itemize}
  \item une seule ligne, avec moins de 200 pare-brise par heure ;
  \item une imprimante CIJ Markem-Imaje 9450c ;
  \item un signal \texttt{line\_demanding} indiquant qu'une impression est
        attendue ;
  \item une caméra de contrôle fournissant un résultat et un grade ;
  \item une durée d'utilisation de 90 jours après ouverture pour les fluides.
\end{itemize}

\begin{warningbox}
Ces valeurs servent au développement. Le modèle exact de l'imprimante, la
liste des tags et la durée réelle après ouverture devront être confirmés avec
la documentation du site avant la mise en service.
\end{warningbox}

\section{Choix d'un simulateur}

Attendre l'accès au Litmus Edge aurait bloqué toutes les étapes suivantes. Le
simulateur n'est donc pas un simple générateur de nombres aléatoires. Il
reproduit les situations qui peuvent tromper la supervision : démarrage,
préchauffage, production, attente sans commande, jet laissé actif, défaut,
remplissage, dérive de viscosité et perte temporaire de communication.

Il accélère le temps par un facteur configurable. Une minute réelle peut
représenter une heure simulée. Les tendances et les alertes deviennent ainsi
visibles pendant une démonstration, tout en gardant des horodatages réels dans
la base.
